\documentclass[11pt]{article}

\usepackage[a4paper,margin=1in]{geometry}
\usepackage{amsmath,amssymb,amsthm,mathtools,bm}
\usepackage[T1]{fontenc}
\usepackage{lmodern}
\usepackage{microtype}
\usepackage{hyperref}

\allowdisplaybreaks
\numberwithin{equation}{section}

\newtheorem{theorem}{Theorem}[section]
\newtheorem{proposition}[theorem]{Proposition}
\newtheorem{lemma}[theorem]{Lemma}
\newtheorem{corollary}[theorem]{Corollary}
\newtheorem{remark}[theorem]{Remark}
\newtheorem{definition}[theorem]{Definition}

\title{Simply Connected Pseudospectral Components for the\\ Open-Boundary Hatano--Nelson Matrix}
\author{}
\date{}

\begin{document}
\maketitle

\begin{abstract}
For the real nonnormal tridiagonal Toeplitz matrix
\[
\mathbf A_n(a)=\operatorname{tridiag}(a,0,1),\qquad n>1,\quad 0<a<1,
\]
we prove that every connected component of every pseudospectrum
\[
\Sigma_\varepsilon=\sigma_\varepsilon\bigl(\mathbf A_n(a)\bigr)
\]
is simply connected.  The key analytic statement is that, for every fixed
$x\in\mathbb{R}$, the least singular value
\[
s(x,y)=s_{\min}\!\bigl((x+iy)\mathbf I_n-\mathbf A_n(a)\bigr)
\]
is nondecreasing for $y\ge 0$.  The proof combines a topological reduction,
a singular-vector current identity, a Perron--Frobenius analysis on the
imaginary axis, continuation of phase inequalities in the $x$-variable, and a
no-zero theorem for ground-state eigenvectors obtained from one-sided Schur
complements of
\[
\mathbf H_x(y)=\bigl((x+iy)\mathbf I_n-\mathbf A_n(a)\bigr)^*
\bigl((x+iy)\mathbf I_n-\mathbf A_n(a)\bigr).
\]
\end{abstract}

\section{Introduction and main results}

Throughout, let
\[
\mathbf A_n(a)=
\begin{bmatrix}
0 & 1 \\
a & 0 & \ddots \\
& \ddots & \ddots & \ddots \\
&& \ddots & 0 & 1 \\
&&& a & 0
\end{bmatrix}\in\mathbb{R}^{n\times n},
\qquad n>1,\quad 0<a<1.
\]
For $\varepsilon>0$, its $\varepsilon$-pseudospectrum is
\[
\Sigma_\varepsilon
=
\sigma_\varepsilon\bigl(\mathbf A_n(a)\bigr)
=
\left\{
z\in\mathbb{C}:\;
s_{\min}\!\bigl(z\mathbf I_n-\mathbf A_n(a)\bigr)<\varepsilon
\right\}.
\]
If $z=x+iy$, we write
\[
s(x,y)=s_{\min}\!\bigl((x+iy)\mathbf I_n-\mathbf A_n(a)\bigr),
\qquad
\lambda_1(x,y)=s(x,y)^2.
\]

Our main result is the following.

\begin{theorem}[Main theorem]\label{thm:main}
For every $\varepsilon>0$, every connected component of
\[
\Sigma_\varepsilon=\sigma_\varepsilon\bigl(\mathbf A_n(a)\bigr)
\]
is simply connected.
\end{theorem}

The proof is reduced to the following monotonicity statement.

\begin{theorem}[Vertical monotonicity]\label{thm:vertical}
For every fixed $x\in\mathbb{R}$, the function
\[
[0,\infty)\ni y\longmapsto s(x,y)
\]
is nondecreasing.  Equivalently,
\[
\partial_y\lambda_1(x,y)\ge 0,
\qquad x\in\mathbb{R},\ y>0.
\]
\end{theorem}

Section~\ref{sec:topology} proves that Theorem~\ref{thm:main} follows from
Theorem~\ref{thm:vertical}.  The remaining sections establish
Theorem~\ref{thm:vertical}.

\section{Topological reduction}\label{sec:topology}

We begin with the spectrum of $\mathbf A_n(a)$.

\begin{lemma}\label{lem:real-spectrum}
Let
\[
\mathbf T_a=
\begin{bmatrix}
1 \\
& \sqrt a \\
&& a \\
&&& \ddots \\
&&&& a^{(n-1)/2}
\end{bmatrix}.
\]
Then
\[
\mathbf T_a^{-1}\mathbf A_n(a)\mathbf T_a
=
\sqrt a\,
\begin{bmatrix}
0 & 1 \\
1 & 0 & \ddots \\
& \ddots & \ddots & \ddots \\
&& \ddots & 0 & 1 \\
&&& 1 & 0
\end{bmatrix}.
\]
Hence the eigenvalues of $\mathbf A_n(a)$ are real and simple:
\[
\sigma\bigl(\mathbf A_n(a)\bigr)
=
\left\{
2\sqrt a\cos\!\left(\frac{k\pi}{n+1}\right):
k=1,\dots,n
\right\}.
\]
\end{lemma}

\begin{proof}
A direct computation gives
\[
(\mathbf T_a^{-1}\mathbf A_n(a)\mathbf T_a)_{j,j+1}
=
a^{-(j-1)/2}a^{j/2}=\sqrt a
\]
and
\[
(\mathbf T_a^{-1}\mathbf A_n(a)\mathbf T_a)_{j+1,j}
=
a^{-j/2}a\,a^{(j-1)/2}=\sqrt a.
\]
Thus $\mathbf T_a^{-1}\mathbf A_n(a)\mathbf T_a=\sqrt a\,\operatorname{tridiag}(1,0,1)$, whose
Dirichlet eigenvalues are the stated numbers.
\end{proof}

\begin{lemma}\label{lem:real-axis}
Every connected component of $\Sigma_\varepsilon$ meets the real axis.
\end{lemma}

\begin{proof}
Let $\Omega$ be a connected component of $\Sigma_\varepsilon$.  Since
\[
s_{\min}\!\bigl(z\mathbf I_n-\mathbf A_n(a)\bigr)\ge |z|-\|\mathbf A_n(a)\|_2,
\]
the set $\Sigma_\varepsilon$ is bounded, hence so is $\Omega$.

Assume that $\Omega\cap\sigma(\mathbf A_n(a))=\varnothing$.  Then
\[
\varphi(z)=\bigl\|(z\mathbf I_n-\mathbf A_n(a))^{-1}\bigr\|_2
\]
is continuous on $\overline\Omega$, subharmonic on $\Omega$, and satisfies
$\varphi(z)>\varepsilon^{-1}$ for $z\in\Omega$.
Because $\Omega$ is a connected component of the open set $\Sigma_\varepsilon$,
we have $\partial\Omega\subset\partial\Sigma_\varepsilon$.  Since
\[
\Sigma_\varepsilon=\{z\in\mathbb{C}:s(z)<\varepsilon\}
\]
and $s$ is continuous, one has $s=\varepsilon$ on $\partial\Sigma_\varepsilon$, hence
\[
\varphi(z)=\varepsilon^{-1},\qquad z\in\partial\Omega.
\]
This contradicts the maximum principle for subharmonic functions on bounded
domains.  Therefore $\Omega$ contains an eigenvalue of $\mathbf A_n(a)$, and
Lemma~\ref{lem:real-spectrum} implies $\Omega\cap\mathbb{R}\neq\varnothing$.
\end{proof}

\begin{lemma}\label{lem:conj}
For all $x,y\in\mathbb{R}$,
\[
s(x,-y)=s(x,y).
\]
Hence $\Sigma_\varepsilon$ is invariant under complex conjugation.
\end{lemma}

\begin{proof}
Since $\mathbf A_n(a)$ has real entries,
\[
\overline{(x+iy)\mathbf I_n-\mathbf A_n(a)}
=
(x-iy)\mathbf I_n-\mathbf A_n(a),
\]
and complex conjugation preserves singular values.
\end{proof}

\begin{proposition}\label{prop:topological-reduction}
Assume Theorem~\ref{thm:vertical}.  Then for every $x\in\mathbb{R}$, the vertical
section
\[
\Sigma_\varepsilon\cap\bigl(\{x\}+i\mathbb{R}\bigr)
\]
is either empty or of the form
\[
\left\{x+iy:\ |y|<h_\varepsilon(x)\right\}
\]
for some $h_\varepsilon(x)>0$.  Consequently, every connected component $\Omega$
of $\Sigma_\varepsilon$ is of the form
\[
\Omega
=
\left\{
x+iy:\ x\in I_\Omega,\ |y|<h_\Omega(x)
\right\},
\]
where $I_\Omega\subset\mathbb{R}$ is an open interval and $h_\Omega:I_\Omega\to(0,\infty)$.
In particular, $\Omega$ is contractible.
\end{proposition}

\begin{proof}
Fix $x\in\mathbb{R}$ and set
\[
V_x=\{y\in\mathbb{R}:\ x+iy\in\Sigma_\varepsilon\}.
\]
By Lemma~\ref{lem:conj}, the set $V_x$ is symmetric under $y\mapsto -y$.
By Theorem~\ref{thm:vertical}, the set
\[
\{y\ge 0:\ s(x,y)<\varepsilon\}
\]
is either empty or an interval of the form $[0,h_\varepsilon(x))$.
Hence $V_x$ is either empty or $(-h_\varepsilon(x),h_\varepsilon(x))$.

Now let $\Omega$ be a connected component of $\Sigma_\varepsilon$, and define
\[
I_\Omega
=
\{x\in\mathbb{R}:\ \Omega\cap(\{x\}+i\mathbb{R})\neq\varnothing\}.
\]
Since the projection $x+iy\mapsto x$ is continuous and $\Omega$ is connected,
$I_\Omega$ is connected; since $\Omega$ is open, $I_\Omega$ is open.  Thus
$I_\Omega$ is an open interval.

By Lemma~\ref{lem:real-axis}, $\Omega$ contains a real point.  Since complex
conjugation preserves $\Sigma_\varepsilon$ and fixes that real point, it preserves
$\Omega$.  Therefore every fiber
\[
\Omega_x=\{y\in\mathbb{R}:\ x+iy\in\Omega\},\qquad x\in I_\Omega,
\]
is symmetric about $0$.  Because $V_x$ is connected and is partitioned by
the open sets corresponding to the component slices, $\Omega_x$ must equal
the whole of $V_x$.  Thus
\[
\Omega_x=(-h_\Omega(x),h_\Omega(x))
\]
for some $h_\Omega(x)>0$, and the stated representation follows.

Finally,
\[
H:[0,1]\times \Omega\longrightarrow \Omega,
\qquad
H\bigl(t,x+iy\bigr)=x+i(1-t)y,
\]
is a deformation retraction of $\Omega$ onto $\Omega\cap\mathbb{R}=I_\Omega$.
Hence $\Omega$ is contractible.
\end{proof}

\section{Singular vectors, currents, and continuation}\label{sec:singular}

For $y>0$ set
\[
\mathbf B_{x,y}=(x+iy)\mathbf I_n-\mathbf A_n(a),
\qquad
\mathbf H_x(y)=\mathbf B_{x,y}^*\mathbf B_{x,y}.
\]
We also introduce the reversal matrix $\bm{\Pi}\in\mathbb{R}^{n\times n}$ by
\[
\bm{\Pi}_{j,n+1-j}=1,\qquad \bm{\Pi}_{jk}=0\ \text{otherwise}.
\]

\begin{lemma}[Canonical gauge]\label{lem:canonical}
One has
\[
\bm{\Pi}\,\mathbf A_n(a)\,\bm{\Pi}=\mathbf A_n(a)^{\mathsf T}.
\]
Suppose one of the variables $x$ or $y$ varies on an interval on which
$\lambda_1(x,y)$ is simple.  Then there exist real-analytic unit vectors
$\bm u(x,y),\bm v(x,y)\in\mathbb{C}^n$ such that
\[
\mathbf B_{x,y}\bm v=s(x,y)\bm u,
\qquad
\mathbf B_{x,y}^*\bm u=s(x,y)\bm v,
\qquad
\bm u=\bm{\Pi}\,\overline{\bm v}.
\]
\end{lemma}

\begin{proof}
The identity $\bm{\Pi}\mathbf A_n(a)\bm{\Pi}=\mathbf A_n(a)^{\mathsf T}$ is
immediate from the Toeplitz form of $\mathbf A_n(a)$.

On an interval of simplicity, standard perturbation theory yields a
real-analytic unit eigenvector $\bm v$ of $\mathbf H_x(y)$ for the simple
eigenvalue $\lambda_1(x,y)=s(x,y)^2$.  Since $y>0$ and the spectrum of
$\mathbf A_n(a)$ is real, $s(x,y)>0$.  Hence
\[
\bm u=s(x,y)^{-1}\mathbf B_{x,y}\bm v
\]
is real-analytic and satisfies
\[
\mathbf B_{x,y}\bm v=s(x,y)\bm u,
\qquad
\mathbf B_{x,y}^*\bm u=s(x,y)\bm v.
\]

Because $\bm{\Pi}\mathbf B_{x,y}=\mathbf B_{x,y}^{\mathsf T}\bm{\Pi}$,
the pair
\[
\bm{\Pi}\,\overline{\bm u},
\qquad
\bm{\Pi}\,\overline{\bm v}
\]
is another singular pair for the same simple singular value.  Therefore there
exists a real-analytic unimodular scalar $\alpha$ such that
\[
\bm{\Pi}\,\overline{\bm u}=\alpha\,\bm v,
\qquad
\bm{\Pi}\,\overline{\bm v}=\alpha\,\bm u.
\]
On an interval one may choose a real-analytic unimodular square root
$\beta$ of $\alpha$.  Replacing $(\bm u,\bm v)$ by
$(\beta\bm u,\beta\bm v)$ yields
\[
\bm u=\bm{\Pi}\,\overline{\bm v}.
\]
\end{proof}

\begin{proposition}[Current identity]\label{prop:current}
Fix $x\in\mathbb{R}$, and let $I\subset(0,\infty)$ be an interval on which
$\lambda_1(x,y)$ is simple.  Let $\bm v(y)$ be the canonical right singular
vector from Lemma~\ref{lem:canonical}, written as
\[
\bm v(y)=
\begin{bmatrix}
v_1(y)\\
\vdots\\
v_n(y)
\end{bmatrix},
\qquad
v_0(y)=v_{n+1}(y)=0.
\]
Define
\[
J_j(y)=\Im\!\bigl(\overline{v_j(y)}\,v_{j+1}(y)\bigr),
\qquad
j=1,\dots,n-1,
\]
and set $J_0(y)=J_n(y)=0$.  Then for $j=1,\dots,n$,
\begin{equation}\label{eq:current}
J_j-aJ_{j-1}
=
y|v_j|^2+s(x,y)\,\Im\!\bigl(v_jv_{n+1-j}\bigr),
\end{equation}
and
\begin{equation}\label{eq:lambda-derivative}
\frac{d}{dy}\lambda_1(x,y)
=
2y-2(1-a)\sum_{j=1}^{n-1}J_j(y).
\end{equation}
\end{proposition}

\begin{proof}
Let
\[
\mathbf B_y=(x+iy)\mathbf I_n-\mathbf A_n(a),
\qquad
s=s(x,y).
\]
Since
\[
\mathbf B_y\bm v=s\bm u,
\qquad
\bm u=\bm{\Pi}\,\overline{\bm v},
\]
the $j$th component reads
\[
(x+iy)v_j-av_{j-1}-v_{j+1}
=
s\,\overline{v_{n+1-j}}.
\]
Multiplying by $\overline{v_j}$ and taking imaginary parts gives
\[
y|v_j|^2-a\,\Im(\overline{v_j}v_{j-1})-\Im(\overline{v_j}v_{j+1})
=
-s\,\Im(v_jv_{n+1-j}),
\]
which is exactly \eqref{eq:current}.

For the derivative, the standard singular-value perturbation formula gives
\[
\frac{d}{dy}s(x,y)
=
\Re\!\bigl(\bm u^* i\bm v\bigr)
=
-\Im(\bm u^*\bm v).
\]
Since $\bm u=\bm{\Pi}\,\overline{\bm v}$,
\[
\bm u^*\bm v
=
\bm v^{\mathsf T}\bm{\Pi}\bm v
=
\sum_{j=1}^n v_jv_{n+1-j}.
\]
Summing \eqref{eq:current} over $j=1,\dots,n$ and using $J_0=J_n=0$ and
$\|\bm v\|_2=1$ yields
\[
(1-a)\sum_{j=1}^{n-1}J_j
=
y+s\,\Im(\bm u^*\bm v).
\]
Therefore
\[
s\,\frac{d}{dy}s
=
y-(1-a)\sum_{j=1}^{n-1}J_j.
\]
Multiplying by $2$ gives \eqref{eq:lambda-derivative}.
\end{proof}

\begin{lemma}[A real telescoping invariant]\label{lem:W}
Assume $\lambda_1(x,y)$ is simple and let $\bm v$ be in canonical gauge.
Define
\[
W_j=v_{j+1}v_{n+1-j}-a\,v_jv_{n-j},
\qquad
j=0,\dots,n,
\]
with $v_0=v_{n+1}=0$.  Then
\[
W_0=0,
\qquad
W_j-W_{j-1}
=
s(x,y)\bigl(|v_j|^2-|v_{n+1-j}|^2\bigr)\in\mathbb{R},
\]
so that
\[
W_j\in\mathbb{R},\qquad j=0,\dots,n.
\]
\end{lemma}

\begin{proof}
Using the singular-vector equations at the indices $j$ and $n+1-j$,
\[
v_{j+1}+av_{j-1}=(x+iy)v_j-s\,\overline{v_{n+1-j}},
\]
\[
v_{n+2-j}+av_{n-j}=(x+iy)v_{n+1-j}-s\,\overline{v_j},
\]
we obtain
\begin{align*}
W_j-W_{j-1}
&=
v_{n+1-j}\bigl(v_{j+1}+av_{j-1}\bigr)
-
v_j\bigl(v_{n+2-j}+av_{n-j}\bigr) \\
&=
s\bigl(|v_j|^2-|v_{n+1-j}|^2\bigr)\in\mathbb{R}.
\end{align*}
Since $W_0=0$, the claim follows.
\end{proof}

\begin{lemma}[Boundary entries on simple branches]\label{lem:boundary}
Assume $\lambda_1(x,y)$ is simple and $y>0$.  In canonical gauge,
\[
v_1\neq0,
\qquad
v_n\neq0.
\]
\end{lemma}

\begin{proof}
If $n=2$, the claim is immediate from the singular-vector equations because
$x+iy\neq0$.

Assume first that $v_n=0$.  If also $v_1=0$, then the equations at $j=1$
and $j=n$ give $v_2=0$ and $v_{n-1}=0$, and repeated propagation yields
$\bm v=0$, impossible.  Hence $v_1\neq0$.

From the equations at $j=1$ and $j=n-1$ we get
\[
v_2=(x+iy)v_1
\]
and
\[
(x+iy)v_{n-1}-av_{n-2}=s(x-iy)\overline{v_1}.
\]
Therefore
\[
W_{n-2}
=
v_{n-1}v_2-av_{n-2}v_1
=
s(x-iy)|v_1|^2.
\]
By Lemma~\ref{lem:W}, the number $W_{n-2}$ is real, which is impossible
because $y>0$ and $v_1\neq0$.  Thus $v_n\neq0$.

The proof that $v_1\neq0$ is analogous.  If $v_1=0$, then
\[
v_2=-s\,\overline{v_n},
\qquad
av_{n-1}=(x+iy)v_n.
\]
Using the equation at $j=2$, one obtains
\[
W_2
=
-\frac{s}{a}(x-iy)|v_n|^2,
\]
again contradicting Lemma~\ref{lem:W}.
\end{proof}

\begin{proposition}[Imaginary-axis seed]\label{prop:x0}
For every $y>0$, the smallest eigenvalue $\lambda_1(0,y)$ of $\mathbf H_0(y)$
is simple.  In canonical gauge, the corresponding right singular vector can be
written as
\[
\bm v(0,y)
=
\beta\,
\begin{bmatrix}
w_1\\
iw_2\\
\vdots\\
i^{\,n-1}w_n
\end{bmatrix},
\qquad
w_j>0,
\]
for some unimodular scalar $\beta$.  Consequently,
\[
\Im\!\left(\frac{v_{j+1}(0,y)}{v_j(0,y)}\right)>0,
\qquad
j=1,\dots,n-1,
\]
and
\[
\Im\!\left(\frac{\overline{v_{n+1-j}(0,y)}}{v_j(0,y)}\right)>0,
\qquad
j=1,\dots,n.
\]
\end{proposition}

\begin{proof}
Set
\[
\mathbf U=
\begin{bmatrix}
1 \\
& i \\
&& i^2 \\
&&& \ddots \\
&&&& i^{\,n-1}
\end{bmatrix},
\qquad
\widehat{\mathbf H}_y=\mathbf U^* \mathbf H_0(y)\mathbf U.
\]
A direct computation shows that $\widehat{\mathbf H}_y$ is real symmetric,
with diagonal entries
\[
a^2+y^2,\quad 1+a^2+y^2,\dots,1+a^2+y^2,\quad 1+y^2,
\]
first off-diagonal entries
\[
-(1-a)y<0,
\]
and second off-diagonal entries
\[
-a<0.
\]
Hence $\widehat{\mathbf H}_y$ is a real symmetric irreducible $Z$-matrix.
Choosing $\Lambda>\max_j (\widehat{\mathbf H}_y)_{jj}$, the matrix
$\Lambda\mathbf I_n-\widehat{\mathbf H}_y$ is entrywise nonnegative and
irreducible.  By Perron--Frobenius, the smallest eigenvalue of
$\widehat{\mathbf H}_y$ is simple and has a strictly positive eigenvector
\[
\bm w=
\begin{bmatrix}
w_1\\
\vdots\\
w_n
\end{bmatrix},
\qquad
w_j>0.
\]
Therefore the right singular vector of $\mathbf B_{0,y}=iy\mathbf I_n-\mathbf A_n(a)$
for $s(0,y)$ can be written as
\[
v_j=\beta\,i^{\,j-1}w_j,
\qquad
j=1,\dots,n,
\]
for some unimodular $\beta$.

Insert this into
\[
iy\,v_j-av_{j-1}-v_{j+1}=s(0,y)\,\overline{v_{n+1-j}}.
\]
After dividing by $\beta i^j$ we get
\[
w_{j+1}-yw_j-aw_{j-1}
=
-\gamma\,s(0,y)\,w_{n+1-j},
\qquad
\gamma=\frac{\overline\beta}{\beta}(-i)^n.
\]
At $j=n$, this yields
\[
yw_n+aw_{n-1}=\gamma\,s(0,y)\,w_1.
\]
The left-hand side is strictly positive, hence $\gamma=1$.

Now
\[
\frac{v_{j+1}(0,y)}{v_j(0,y)}
=
i\,\frac{w_{j+1}}{w_j},
\]
and therefore the first set of inequalities holds.  Similarly,
\[
\frac{\overline{v_{n+1-j}(0,y)}}{v_j(0,y)}
=
i\,\frac{w_{n+1-j}}{w_j},
\]
which proves the second.
\end{proof}

\begin{lemma}[Angular identity]\label{lem:angular}
Assume $\lambda_1(x,y)$ is simple and let $\bm v$ be in canonical gauge.
For $j=1,\dots,n-1$, define
\[
r_j=\frac{v_{j+1}}{v_j},
\qquad
\rho_j=r_{n-j},
\qquad
p_j=\frac{\overline{v_{n+1-j}}}{v_j},
\qquad
r_0=\infty,
\]
and
\[
\alpha_j=(x+iy)-\frac{a}{r_{j-1}},
\qquad
\beta_j=(x-iy)-\overline{\rho_{j-1}}.
\]
Then
\[
\Im\!\bigl(\alpha_j\overline{p_j}\bigr)
=
-\Im\!\bigl(\beta_j p_j\bigr).
\]
\end{lemma}

\begin{proof}
From the singular-vector equation at index $j$,
\[
\left((x+iy)-\frac{a}{r_{j-1}}\right)v_j
=
v_{j+1}+s\,\overline{v_{n+1-j}},
\]
we obtain
\[
\alpha_j\overline{p_j}
=
\frac{v_{j+1}v_{n+1-j}+s|v_{n+1-j}|^2}{|v_j|^2},
\]
hence
\[
\Im\!\bigl(\alpha_j\overline{p_j}\bigr)
=
\frac{\Im(v_{j+1}v_{n+1-j})}{|v_j|^2}.
\]

Conjugating the singular-vector equation at index $n+1-j$ yields
\[
(x-iy)\overline{v_{n+1-j}}
-
a\overline{v_{n-j}}
-
\overline{v_{n+2-j}}
=
sv_j.
\]
Therefore
\[
\beta_j p_j
=
s+a\,\frac{\overline{v_{n-j}}}{v_j},
\]
and thus
\[
\Im(\beta_jp_j)
=
-\frac{a\,\Im(v_jv_{n-j})}{|v_j|^2}.
\]

Finally, Lemma~\ref{lem:W} gives
\[
\Im(v_{j+1}v_{n+1-j})=a\,\Im(v_jv_{n-j}),
\]
and the claim follows.
\end{proof}

\begin{proposition}[Phase continuation and branchwise monotonicity]
\label{prop:continuation}
Fix $y>0$, and let $I\subset\mathbb{R}$ be an open interval on which
$\lambda_1(x,y)$ is simple.  Let
\[
\bm v(x)=
\begin{bmatrix}
v_1(x)\\
\vdots\\
v_n(x)
\end{bmatrix}
\]
be a real-analytic canonical right singular vector, and let $J$ be a connected
component of
\[
\left\{
x\in I:\ v_1(x)\cdots v_n(x)\neq0
\right\}.
\]
Assume that for some $x_0\in J$,
\[
\Im\!\left(\frac{v_{j+1}(x_0)}{v_j(x_0)}\right)>0,
\qquad
j=1,\dots,n-1,
\]
and
\[
\Im\!\left(\frac{\overline{v_{n+1-j}(x_0)}}{v_j(x_0)}\right)>0,
\qquad
j=1,\dots,n.
\]
Then the same inequalities hold for every $x\in J$.  In particular,
\[
\partial_y\lambda_1(x,y)\ge0,
\qquad x\in J.
\]
\end{proposition}

\begin{proof}
\emph{Step 1: continuation of the ratios $r_j=v_{j+1}/v_j$.}
Since $v_j(x)\neq0$ on $J$, we may write
\[
v_j(x)=t_j(x)e^{i\phi_j(x)},
\qquad
t_j(x)>0,
\]
with continuous arguments $\phi_j$.  Set
\[
\delta_j(x)=\phi_{j+1}(x)-\phi_j(x),
\qquad
j=1,\dots,n-1.
\]
Then
\[
\frac{v_{j+1}(x)}{v_j(x)}
=
\frac{t_{j+1}(x)}{t_j(x)}e^{i\delta_j(x)},
\]
so $\Im(v_{j+1}/v_j)>0$ is equivalent to $\delta_j\in(0,\pi)$.

Write
\[
c=(1+a)x-i(1-a)y=|c|e^{-i\vartheta},
\qquad
\vartheta\in(0,\pi).
\]
A direct computation gives
\[
\langle \bm v,\mathbf H_x(y)\bm v\rangle
=
\sum_{j=1}^n \gamma_j t_j^2+\mathcal P_x(\delta_1,\dots,\delta_{n-1}),
\]
where $\gamma_j\in\mathbb{R}$ and
\[
\mathcal P_x
=
-2|c|\sum_{j=1}^{n-1}t_jt_{j+1}\cos(\delta_j-\vartheta)
+
2a\sum_{j=1}^{n-2}t_jt_{j+2}\cos(\delta_j+\delta_{j+1}).
\]
For fixed amplitudes $t_1,\dots,t_n$, the phase vector
\[
(\delta_1(x),\dots,\delta_{n-1}(x))
\]
minimizes $\mathcal P_x$, because $\bm v(x)$ is a ground state of
$\mathbf H_x(y)$.

Let
\[
E=
\left\{
x\in J:\ \delta_j(x)\in(0,\pi)\ \text{for all }j=1,\dots,n-1
\right\}.
\]
By assumption, $x_0\in E$, so $E\neq\varnothing$.  Clearly $E$ is open in $J$.
We show that $E$ is closed in $J$.

Take $x_*\in\overline E\cap J$.  Then $\delta_j(x_*)\in[0,\pi]$ for all $j$.
If every $\delta_j(x_*)$ lies in $(0,\pi)$, then $x_*\in E$.  Thus suppose
$\delta_k(x_*)\in\{0,\pi\}$ for some $k$.

Vary the tail phases by
\[
\phi_m\mapsto \phi_m+\varepsilon,\qquad m\ge k+1.
\]
This changes only $\delta_k$, replacing it by $\delta_k+\varepsilon$.  Therefore the
derivative at $\varepsilon=0$ of the phase functional equals
\begin{align*}
f_k'(0)
&=
2|c|\,t_kt_{k+1}\sin(\delta_k-\vartheta) \\
&\quad
-2a\,\mathbf 1_{\{k\ge2\}}\,t_{k-1}t_{k+1}\sin(\delta_{k-1}+\delta_k) \\
&\quad
-2a\,\mathbf 1_{\{k\le n-2\}}\,t_kt_{k+2}\sin(\delta_k+\delta_{k+1}).
\end{align*}
If $\delta_k(x_*)=0$, then the first term is strictly negative and the other
two are nonpositive, so $f_k'(0)<0$, contradicting minimality.  If
$\delta_k(x_*)=\pi$, then the first term is strictly positive and the other
two are nonnegative, so $f_k'(0)>0$, again contradicting minimality.
Hence no $\delta_k$ can hit $0$ or $\pi$, and therefore $x_*\in E$.
Thus $E$ is closed in $J$, and since $J$ is connected, $E=J$.

We have proved
\[
\Im\!\left(\frac{v_{j+1}(x)}{v_j(x)}\right)>0,
\qquad
j=1,\dots,n-1,\ \ x\in J.
\]
In particular,
\[
\Im \rho_j(x)>0,
\qquad
j=1,\dots,n-1,\ \ x\in J.
\]

\medskip
\noindent
\emph{Step 2: continuation of the reflected phases
$p_j=\overline{v_{n+1-j}}/v_j$.}
For $j=1,\dots,n-1$, suppose that $\Im p_j$ vanishes somewhere on $J$.
By continuity there exists $x_*\in J$ with $p_j(x_*)\in\mathbb{R}$.  Since all
components of $\bm v$ are nonzero on $J$, we have $p_j(x_*)\neq0$.

By Step~1,
\[
\Im\alpha_j(x_*)=
y+\Im\!\left(-\frac{a}{r_{j-1}(x_*)}\right)>0.
\]
Lemma~\ref{lem:angular} therefore gives
\[
\Im\beta_j(x_*)=-\Im\alpha_j(x_*)<0.
\]
Conjugating the singular-vector equation at index $n+1-j$ shows that
\[
\overline{\rho_j}
=
\frac{a}{\beta_j-s/p_j}.
\]
Since $s/p_j\in\mathbb{R}$ at $x_*$, the denominator has negative imaginary part,
and hence $\Im\overline{\rho_j}(x_*)>0$, that is,
\[
\Im\rho_j(x_*)<0,
\]
contradicting Step~1.  Thus
\[
\Im p_j(x)>0,
\qquad
j=1,\dots,n-1,\ \ x\in J.
\]
For $j=n$ we use
\[
p_n(x)=\frac{1}{\overline{p_1(x)}},
\]
which shows $\Im p_n(x)>0$ as well.

\medskip
\noindent
\emph{Step 3: monotonicity in $y$.}
Since
\[
p_j
=
\frac{\overline{v_{n+1-j}}}{v_j}
=
\frac{\overline{v_jv_{n+1-j}}}{|v_j|^2},
\]
we have
\[
\Im p_j
=
-\frac{\Im(v_jv_{n+1-j})}{|v_j|^2}.
\]
Therefore $\Im p_j>0$ implies
\[
\Im(v_jv_{n+1-j})<0.
\]
Using \eqref{eq:current}, we obtain
\[
J_j-aJ_{j-1}\le y|v_j|^2,
\qquad
j=1,\dots,n.
\]
Inductively,
\[
J_j
\le
y\sum_{k=1}^j a^{\,j-k}|v_k|^2,
\qquad
j=1,\dots,n-1.
\]
Summing this inequality gives
\[
(1-a)\sum_{j=1}^{n-1}J_j
\le
y\sum_{k=1}^{n-1}(1-a^{\,n-k})|v_k|^2
\le y.
\]
Now \eqref{eq:lambda-derivative} yields
\[
\partial_y\lambda_1(x,y)\ge0,
\qquad x\in J.
\]
\end{proof}

\section{One-sided positivity and edge Schur complements}\label{sec:schur}

Fix $x\in\mathbb{R}$ and $y>0$, and write
\[
z=x+iy,
\qquad
\lambda=\lambda_1(x,y),
\qquad
c=\overline z+az=(1+a)x-i(1-a)y.
\]
Set
\[
d_-=a^2+|z|^2-\lambda,
\qquad
d=1+a^2+|z|^2-\lambda,
\qquad
d_+=1+|z|^2-\lambda.
\]
Then
\[
\mathbf H_x(y)-\lambda\mathbf I_n
=
\begin{bmatrix}
d_- & -c & a \\
-\overline c & d & -c & a \\
a & -\overline c & d & -c & a \\
& \ddots & \ddots & \ddots & \ddots & \ddots \\
&& a & -\overline c & d & -c \\
&&& a & -\overline c & d_+
\end{bmatrix}.
\]
For an index set $I\subset\{1,\dots,n\}$, we write $\mathbf M[I]$ for the
principal submatrix of $\mathbf M$ indexed by $I$.

For $j=2,\dots,n-1$, define the one-sided truncations
\[
\mathbf L_j
=
\bigl(\mathbf H_x(y)-\lambda\mathbf I_n\bigr)[\{1,\dots,j-1\}],
\qquad
\mathbf R_j
=
\bigl(\mathbf H_x(y)-\lambda\mathbf I_n\bigr)[\{j+1,\dots,n\}].
\]

\begin{proposition}[One-sided positivity]\label{prop:onesided}
For every $j=2,\dots,n-1$,
\[
\mathbf L_j\succ0,
\qquad
\mathbf R_j\succ0.
\]
In particular,
\[
d_->0,
\qquad
d_+>0.
\]
\end{proposition}

\begin{proof}
Since $\mathbf H_x(y)-\lambda\mathbf I_n\succeq0$, every principal submatrix is
positive semidefinite.  It therefore suffices to prove that the kernels of
$\mathbf L_j$ and $\mathbf R_j$ are trivial.

We prove this for $\mathbf R_j$; the proof for $\mathbf L_j$ is symmetric.
Let $\bm w\in\ker\mathbf R_j$, and extend it by zeros to
\[
\widehat{\bm w}=
\begin{bmatrix}
0\\
\vdots\\
0\\
w_1\\
\vdots\\
w_{n-j}
\end{bmatrix}\in\mathbb{C}^n.
\]
Then the rows $j+1,\dots,n$ of
$(\mathbf H_x(y)-\lambda\mathbf I_n)\widehat{\bm w}$ vanish.  Row $j-1$
gives
\[
aw_1=0,
\]
hence $w_1=0$.  If $n-j\ge2$, row $j$ then gives
\[
aw_2=0,
\]
hence $w_2=0$.  The first row of $\mathbf R_j\bm w=0$ now yields
\[
aw_3=0.
\]
Proceeding inductively, all components vanish.  Thus
$\ker\mathbf R_j=\{0\}$, so $\mathbf R_j\succ0$.

The proof for $\mathbf L_j$ is identical.  Finally,
\[
\mathbf L_2=[d_-],
\qquad
\mathbf R_{n-1}=[d_+],
\]
hence $d_->0$ and $d_+>0$.
\end{proof}

\begin{definition}\label{def:edge}
Define
\[
\mathbf T=
\begin{bmatrix}
d & -c\\
-\overline c & d
\end{bmatrix},
\qquad
\mathbf E=
\begin{bmatrix}
a & 0\\
-c & a
\end{bmatrix}.
\]
For $j=3,\dots,n-1$, let $\mathbf S_j^{\mathrm L}$ be the Schur complement of
$\mathbf L_j$ onto its last two coordinates $(j-2,j-1)$.  For
$j=2,\dots,n-2$, let $\mathbf S_j^{\mathrm R}$ be the Schur complement of
$\mathbf R_j$ onto its first two coordinates $(j+1,j+2)$.
\end{definition}

\begin{proposition}[Riccati recursions]\label{prop:schur-rec}
Whenever the indices are meaningful,
\[
\mathbf S_3^{\mathrm L}
=
\begin{bmatrix}
d_- & -c\\
-\overline c & d
\end{bmatrix},
\qquad
\mathbf S_4^{\mathrm L}
=
\begin{bmatrix}
d-\dfrac{|c|^2}{d_-} & -c+\dfrac{a\overline c}{d_-}\\[3mm]
-\overline c+\dfrac{ac}{d_-} & d-\dfrac{a^2}{d_-}
\end{bmatrix},
\]
and
\[
\mathbf S_j^{\mathrm L}
=
\mathbf T-\mathbf E^*(\mathbf S_{j-2}^{\mathrm L})^{-1}\mathbf E,
\qquad
j\ge5.
\]
Similarly,
\[
\mathbf S_{n-2}^{\mathrm R}
=
\begin{bmatrix}
d & -c\\
-\overline c & d_+
\end{bmatrix},
\qquad
\mathbf S_{n-3}^{\mathrm R}
=
\begin{bmatrix}
d-\dfrac{a^2}{d_+} & -c+\dfrac{a\overline c}{d_+}\\[3mm]
-\overline c+\dfrac{ac}{d_+} & d-\dfrac{|c|^2}{d_+}
\end{bmatrix},
\]
and
\[
\mathbf S_j^{\mathrm R}
=
\mathbf T-\mathbf E(\mathbf S_{j+2}^{\mathrm R})^{-1}\mathbf E^*,
\qquad
j\le n-4.
\]
\end{proposition}

\begin{proof}
This is a direct block Schur-complement calculation for the five-diagonal
matrix $\mathbf H_x(y)-\lambda\mathbf I_n$.
\end{proof}

\begin{lemma}[Propagation of the upper half-plane]\label{lem:propagate}
Let
\[
\mathbf S=
\begin{bmatrix}
\alpha & \beta\\
\overline\beta & \delta
\end{bmatrix}\succ0,
\qquad
\Delta=\alpha\delta-|\beta|^2>0.
\]
If
\[
\mathbf T_{\mathrm L}
=
\mathbf T-\mathbf E^*\mathbf S^{-1}\mathbf E
=
\begin{bmatrix}
\alpha' & \beta'\\
\overline{\beta'} & \delta'
\end{bmatrix},
\]
then
\[
\beta'
=
\frac{a^2\beta+a\alpha\,\overline c-c\,\Delta}{\Delta},
\]
and therefore
\[
\Im\beta'
=
\frac{a^2\Im\beta+a\alpha(1-a)y+\Delta(1-a)y}{\Delta}.
\]
If
\[
\mathbf T_{\mathrm R}
=
\mathbf T-\mathbf E\mathbf S^{-1}\mathbf E^*
=
\begin{bmatrix}
\widetilde\alpha & \widetilde\beta\\
\overline{\widetilde\beta} & \widetilde\delta
\end{bmatrix},
\]
then
\[
\widetilde\beta
=
\frac{a^2\beta+a\delta\,\overline c-c\,\Delta}{\Delta},
\]
and therefore
\[
\Im\widetilde\beta
=
\frac{a^2\Im\beta+a\delta(1-a)y+\Delta(1-a)y}{\Delta}.
\]
In particular, if $\Im\beta>0$, then
\[
\Im\beta'>0,
\qquad
\Im\widetilde\beta>0.
\]
\end{lemma}

\begin{proof}
Using
\[
\mathbf S^{-1}
=
\frac{1}{\Delta}
\begin{bmatrix}
\delta & -\beta\\
-\overline\beta & \alpha
\end{bmatrix},
\]
one computes the off-diagonal entries of $\mathbf T_{\mathrm L}$ and
$\mathbf T_{\mathrm R}$ directly.
\end{proof}

\begin{corollary}[Upper-half-plane property]\label{cor:upper}
For every $j=3,\dots,n-1$,
\[
\mathbf S_j^{\mathrm L}
=
\begin{bmatrix}
\alpha_j^{\mathrm L} & \beta_j^{\mathrm L}\\
\overline{\beta_j^{\mathrm L}} & \delta_j^{\mathrm L}
\end{bmatrix}
\quad\Longrightarrow\quad
\Im\beta_j^{\mathrm L}>0.
\]
For every $j=2,\dots,n-2$,
\[
\mathbf S_j^{\mathrm R}
=
\begin{bmatrix}
\alpha_j^{\mathrm R} & \beta_j^{\mathrm R}\\
\overline{\beta_j^{\mathrm R}} & \delta_j^{\mathrm R}
\end{bmatrix}
\quad\Longrightarrow\quad
\Im\beta_j^{\mathrm R}>0.
\]
\end{corollary}

\begin{proof}
The base cases are
\[
\beta_3^{\mathrm L}=-c,
\qquad
\beta_4^{\mathrm L}=-c+\frac{a\overline c}{d_-},
\qquad
\beta_{n-2}^{\mathrm R}=-c,
\qquad
\beta_{n-3}^{\mathrm R}=-c+\frac{a\overline c}{d_+},
\]
all of which have positive imaginary part because $d_\pm>0$.  The recursions
in Proposition~\ref{prop:schur-rec} and Lemma~\ref{lem:propagate} complete
the proof.
\end{proof}

\section{A no-zero theorem for ground states}\label{sec:nozero}

\begin{proposition}[The second boundary layer never vanishes]
\label{prop:second-layer}
Assume $n\ge3$, and let
\[
\bm v=
\begin{bmatrix}
v_1\\
\vdots\\
v_n
\end{bmatrix}
\in
\ker\bigl(\mathbf H_x(y)-\lambda\mathbf I_n\bigr)\setminus\{0\},
\qquad
\lambda=\lambda_1(x,y).
\]
Then
\[
v_2\neq0,
\qquad
v_{n-1}\neq0.
\]
\end{proposition}

\begin{proof}
We prove $v_2\neq0$; the other statement follows by reversing the coordinates.

Assume $v_2=0$.

If $n=3$, then $v_1\neq0$; otherwise $\mathbf R_2[v_3]^{\mathsf T}=0$ would
contradict Proposition~\ref{prop:onesided}.  Row $3$ gives
\[
d_+v_3+av_1=0,
\qquad\text{so}\qquad
v_3=-\frac{a}{d_+}v_1.
\]
Row $2$ then gives
\[
-\overline c\,v_1-c\,v_3=0,
\]
hence
\[
\overline c=\frac{a}{d_+}c,
\]
which is impossible because $a/d_+>0$ while
$\Im(\overline c)=+(1-a)y$ and $\Im(c)=-(1-a)y$.

Assume now that $n\ge4$.  If $v_1=0$, then
\[
\mathbf R_2
\begin{bmatrix}
v_3\\
\vdots\\
v_n
\end{bmatrix}
=
\bm 0,
\]
contradicting Proposition~\ref{prop:onesided}.  Hence $v_1\neq0$.

Let
\[
\mathbf G_2=\mathbf R_2^{-1},
\qquad
g=\langle \bm e_1,\mathbf G_2\bm e_1\rangle>0,
\qquad
h=\langle \bm e_2,\mathbf G_2\bm e_1\rangle.
\]
From
\[
\mathbf R_2
\begin{bmatrix}
v_3\\
\vdots\\
v_n
\end{bmatrix}
+
av_1\bm e_1
=
\bm 0
\]
we obtain
\[
v_3=-ag\,v_1,
\qquad
v_4=-ah\,v_1.
\]
Row $1$ of
$(\mathbf H_x(y)-\lambda\mathbf I_n)\bm v=\bm0$ gives
\[
d_-v_1+av_3=0,
\]
hence
\[
d_-=a^2g.
\]
Row $2$ gives
\[
-\overline c\,v_1-c\,v_3+av_4=0,
\]
that is,
\[
-\overline c+ag\,c-a^2h=0.
\]
Equivalently,
\[
-c+a\frac{h}{g}
=
-\frac{1}{ag}\,\overline c
=
-\frac{a}{d_-}\,\overline c.
\]
Hence the left-hand side has negative imaginary part.

On the other hand, write
\[
\mathbf S_2^{\mathrm R}
=
\begin{bmatrix}
\alpha_2^{\mathrm R} & \beta_2^{\mathrm R}\\
\overline{\beta_2^{\mathrm R}} & \delta_2^{\mathrm R}
\end{bmatrix}.
\]
By the block inversion formula,
\[
\frac{h}{g}
=
-\frac{\overline{\beta_2^{\mathrm R}}}{\delta_2^{\mathrm R}},
\]
and therefore
\[
-c+a\frac{h}{g}
=
-c-a\,\frac{\overline{\beta_2^{\mathrm R}}}{\delta_2^{\mathrm R}}.
\]
By Corollary~\ref{cor:upper}, $\Im\beta_2^{\mathrm R}>0$, so this number has
positive imaginary part, a contradiction.  Thus $v_2\neq0$.
\end{proof}

\begin{proposition}[Defect resonance for an interior zero]\label{prop:defect}
Assume $n\ge5$, and let
\[
\bm v=
\begin{bmatrix}
v_1\\
\vdots\\
v_n
\end{bmatrix}
\in
\ker\bigl(\mathbf H_x(y)-\lambda\mathbf I_n\bigr)\setminus\{0\},
\qquad
\lambda=\lambda_1(x,y).
\]
Fix $j\in\{3,\dots,n-2\}$ and assume $v_j=0$.
Set
\[
\mathbf G_j^{\mathrm L}=\mathbf L_j^{-1},
\qquad
\mathbf G_j^{\mathrm R}=\mathbf R_j^{-1},
\]
and define
\[
g_j^{\mathrm L}
=
\langle \bm e_{j-1},\mathbf G_j^{\mathrm L}\bm e_{j-1}\rangle>0,
\qquad
h_j^{\mathrm L}
=
\langle \bm e_{j-2},\mathbf G_j^{\mathrm L}\bm e_{j-1}\rangle,
\]
where $\bm e_{j-2},\bm e_{j-1}$ are the last two basis vectors of $\mathbb{C}^{j-1}$,
and
\[
g_j^{\mathrm R}
=
\langle \bm e_1,\mathbf G_j^{\mathrm R}\bm e_1\rangle>0,
\qquad
h_j^{\mathrm R}
=
\langle \bm e_2,\mathbf G_j^{\mathrm R}\bm e_1\rangle,
\]
where $\bm e_1,\bm e_2$ are the first two basis vectors of $\mathbb{C}^{n-j}$.
Then:
\begin{enumerate}
\item
\[
v_{j-1}=-ag_j^{\mathrm L}v_{j+1},
\qquad
v_{j+1}=-ag_j^{\mathrm R}v_{j-1}.
\]
In particular,
\[
v_{j-1}\neq0,
\qquad
v_{j+1}\neq0,
\qquad
a^2g_j^{\mathrm L}g_j^{\mathrm R}=1.
\]

\item
If
\[
m_j^{\mathrm L}=a\frac{h_j^{\mathrm L}}{g_j^{\mathrm L}}-\overline c,
\qquad
m_j^{\mathrm R}=-c+a\frac{h_j^{\mathrm R}}{g_j^{\mathrm R}},
\]
then
\[
m_j^{\mathrm L}v_{j-1}+m_j^{\mathrm R}v_{j+1}=0,
\]
hence
\[
m_j^{\mathrm R}=ag_j^{\mathrm L}m_j^{\mathrm L}.
\]
\end{enumerate}
\end{proposition}

\begin{proof}
Split
\[
\bm v=
\begin{bmatrix}
v_1\\
\vdots\\
v_{j-1}\\
0\\
v_{j+1}\\
\vdots\\
v_n
\end{bmatrix}
=
\begin{bmatrix}
\bm u\\
0\\
\bm w
\end{bmatrix}.
\]
The equations on the left and right of the defect read
\[
\mathbf L_j\bm u+av_{j+1}\bm e_{j-1}=\bm0,
\qquad
\mathbf R_j\bm w+av_{j-1}\bm e_1=\bm0.
\]
Therefore
\[
\bm u=-av_{j+1}\mathbf G_j^{\mathrm L}\bm e_{j-1},
\qquad
\bm w=-av_{j-1}\mathbf G_j^{\mathrm R}\bm e_1,
\]
which gives
\[
v_{j-1}=-ag_j^{\mathrm L}v_{j+1},
\qquad
v_{j+1}=-ag_j^{\mathrm R}v_{j-1}.
\]
If $v_{j+1}=0$, then $\bm u=\bm0$, hence $v_{j-1}=0$, and then $\bm w=\bm0$,
contrary to $\bm v\neq0$.  Thus $v_{j\pm1}\neq0$, and
$a^2g_j^{\mathrm L}g_j^{\mathrm R}=1$ follows immediately.

Moreover,
\[
v_{j-2}=\frac{h_j^{\mathrm L}}{g_j^{\mathrm L}}v_{j-1},
\qquad
v_{j+2}=\frac{h_j^{\mathrm R}}{g_j^{\mathrm R}}v_{j+1}.
\]
The $j$th row of
$(\mathbf H_x(y)-\lambda\mathbf I_n)\bm v=\bm0$ is
\[
av_{j-2}-\overline c\,v_{j-1}-c\,v_{j+1}+av_{j+2}=0,
\]
which is equivalent to
\[
m_j^{\mathrm L}v_{j-1}+m_j^{\mathrm R}v_{j+1}=0.
\]
Using $v_{j-1}=-ag_j^{\mathrm L}v_{j+1}$ and $v_{j+1}\neq0$ yields the last
identity.
\end{proof}

\begin{proposition}[Sign separation for the defect Weyl quotients]
\label{prop:sign}
For $j=3,\dots,n-2$, write
\[
\mathbf S_j^{\mathrm L}
=
\begin{bmatrix}
\alpha_j^{\mathrm L} & \beta_j^{\mathrm L}\\
\overline{\beta_j^{\mathrm L}} & \delta_j^{\mathrm L}
\end{bmatrix},
\qquad
\mathbf S_j^{\mathrm R}
=
\begin{bmatrix}
\alpha_j^{\mathrm R} & \beta_j^{\mathrm R}\\
\overline{\beta_j^{\mathrm R}} & \delta_j^{\mathrm R}
\end{bmatrix}.
\]
Then
\[
\frac{h_j^{\mathrm L}}{g_j^{\mathrm L}}
=
-\frac{\beta_j^{\mathrm L}}{\alpha_j^{\mathrm L}},
\qquad
\frac{h_j^{\mathrm R}}{g_j^{\mathrm R}}
=
-\frac{\overline{\beta_j^{\mathrm R}}}{\delta_j^{\mathrm R}},
\]
and therefore
\[
m_j^{\mathrm L}
=
-\overline c-a\,\frac{\beta_j^{\mathrm L}}{\alpha_j^{\mathrm L}},
\qquad
m_j^{\mathrm R}
=
-c-a\,\frac{\overline{\beta_j^{\mathrm R}}}{\delta_j^{\mathrm R}}.
\]
In particular,
\[
\Im m_j^{\mathrm L}<0<\Im m_j^{\mathrm R}.
\]
\end{proposition}

\begin{proof}
Since $\mathbf S_j^{\mathrm L}$ is the Schur complement of $\mathbf L_j$ onto the
last two coordinates, the lower-right $2\times2$ block of $\mathbf L_j^{-1}$ is
$(\mathbf S_j^{\mathrm L})^{-1}$.  Hence
\[
g_j^{\mathrm L}=\bigl((\mathbf S_j^{\mathrm L})^{-1}\bigr)_{22},
\qquad
h_j^{\mathrm L}=\bigl((\mathbf S_j^{\mathrm L})^{-1}\bigr)_{12}.
\]
Since
\[
(\mathbf S_j^{\mathrm L})^{-1}
=
\frac{1}{\alpha_j^{\mathrm L}\delta_j^{\mathrm L}-|\beta_j^{\mathrm L}|^2}
\begin{bmatrix}
\delta_j^{\mathrm L} & -\beta_j^{\mathrm L}\\
-\overline{\beta_j^{\mathrm L}} & \alpha_j^{\mathrm L}
\end{bmatrix},
\]
the first ratio formula follows.  The proof for $\mathbf S_j^{\mathrm R}$ is
identical.

By Corollary~\ref{cor:upper},
\[
\Im\beta_j^{\mathrm L}>0,
\qquad
\Im\beta_j^{\mathrm R}>0.
\]
Since $\alpha_j^{\mathrm L}>0$ and $\delta_j^{\mathrm R}>0$, we get
\[
\Im m_j^{\mathrm L}
=
-(1-a)y-a\,\frac{\Im\beta_j^{\mathrm L}}{\alpha_j^{\mathrm L}}<0
\]
and
\[
\Im m_j^{\mathrm R}
=
(1-a)y+a\,\frac{\Im\beta_j^{\mathrm R}}{\delta_j^{\mathrm R}}>0.
\]
\end{proof}

\begin{theorem}[No interior zeros for a ground state]\label{thm:nozero}
Assume $n\ge3$, and let
\[
\bm v\in
\ker\bigl(\mathbf H_x(y)-\lambda_1(x,y)\mathbf I_n\bigr)\setminus\{0\}.
\]
Then
\[
v_j\neq0,
\qquad
j=2,\dots,n-1.
\]
\end{theorem}

\begin{proof}
By Proposition~\ref{prop:second-layer},
\[
v_2\neq0,
\qquad
v_{n-1}\neq0.
\]
Thus only the indices $j=3,\dots,n-2$ remain.  If such a $j$ satisfied
$v_j=0$, then Proposition~\ref{prop:defect} would imply
\[
m_j^{\mathrm R}=ag_j^{\mathrm L}m_j^{\mathrm L}
\]
with $ag_j^{\mathrm L}>0$.  Hence $m_j^{\mathrm R}$ would be a positive real
multiple of $m_j^{\mathrm L}$, contradicting Proposition~\ref{prop:sign}.
\end{proof}

\begin{corollary}[Global simplicity of the lowest branch]\label{cor:simple}
Assume $n\ge3$.  For every $x\in\mathbb{R}$ and $y>0$, the smallest eigenvalue
\[
\lambda_1(x,y)=\lambda_{\min}\!\bigl(\mathbf H_x(y)\bigr)
\]
is simple.
\end{corollary}

\begin{proof}
If the eigenspace had dimension at least $2$, then the linear map
\[
\ker\bigl(\mathbf H_x(y)-\lambda_1(x,y)\mathbf I_n\bigr)\longrightarrow\mathbb{C},
\qquad
\bm v\longmapsto v_2,
\]
would have a nontrivial kernel.  This would contradict Theorem~\ref{thm:nozero}.
\end{proof}

\section{Proof of vertical monotonicity}\label{sec:proof}

\begin{proposition}[The case $n=2$]\label{prop:n2}
If $n=2$, then for all $x\in\mathbb{R}$ and $y\ge0$,
\[
\partial_y\lambda_1(x,y)\ge0.
\]
\end{proposition}

\begin{proof}
For $n=2$,
\[
\mathbf B_{x,y}
=
\begin{bmatrix}
x+iy & -1\\
-a & x+iy
\end{bmatrix},
\]
and hence
\[
\mathbf H_x(y)
=
\begin{bmatrix}
a^2+x^2+y^2 & -(a+1)x+i(1-a)y\\
-(a+1)x-i(1-a)y & 1+x^2+y^2
\end{bmatrix}.
\]
Its smaller eigenvalue is
\[
\lambda_1(x,y)
=
x^2+y^2+\frac{1+a^2}{2}
-\frac12\sqrt{(1-a^2)^2+4(a+1)^2x^2+4(1-a)^2y^2}.
\]
Differentiating gives
\[
\partial_y\lambda_1(x,y)
=
2y\left(
1-\frac{(1-a)^2}
{\sqrt{(1-a^2)^2+4(a+1)^2x^2+4(1-a)^2y^2}}
\right)\ge0,
\]
because
\[
\sqrt{(1-a^2)^2+4(a+1)^2x^2+4(1-a)^2y^2}\ge 1-a^2\ge (1-a)^2.
\]
\end{proof}

\begin{proof}[Proof of Theorem~\ref{thm:vertical}]
If $n=2$, the claim is Proposition~\ref{prop:n2}.  Assume henceforth that
$n\ge3$.

Fix $y_*>0$.  By Corollary~\ref{cor:simple}, the lowest eigenvalue
\[
x\longmapsto \lambda_1(x,y_*)
\]
is simple for every $x\in\mathbb{R}$.  Since $x\mapsto \mathbf H_x(y_*)$ is a
real-analytic family, Lemma~\ref{lem:canonical} yields a real-analytic
canonical right singular vector
\[
\bm v(x)=
\begin{bmatrix}
v_1(x)\\
\vdots\\
v_n(x)
\end{bmatrix},
\qquad
x\in\mathbb{R}.
\]

By Theorem~\ref{thm:nozero},
\[
v_j(x)\neq0,
\qquad
j=2,\dots,n-1,\ \ x\in\mathbb{R}.
\]
By Lemma~\ref{lem:boundary}, also
\[
v_1(x)\neq0,
\qquad
v_n(x)\neq0,
\qquad
x\in\mathbb{R}.
\]
Thus every component of $\bm v(x)$ is nonzero for every $x\in\mathbb{R}$.

At $x=0$, Proposition~\ref{prop:x0} gives
\[
\Im\!\left(\frac{v_{j+1}(0)}{v_j(0)}\right)>0,
\qquad
j=1,\dots,n-1,
\]
and
\[
\Im\!\left(\frac{\overline{v_{n+1-j}(0)}}{v_j(0)}\right)>0,
\qquad
j=1,\dots,n.
\]
Hence Proposition~\ref{prop:continuation}, applied with $J=\mathbb{R}$, yields
\[
\partial_y\lambda_1(x,y_*)\ge0,
\qquad
x\in\mathbb{R}.
\]
Since $y_*>0$ was arbitrary, we conclude that
\[
\partial_y\lambda_1(x,y)\ge0,
\qquad
x\in\mathbb{R},\ y>0.
\]

Finally, Lemma~\ref{lem:conj} implies that $\lambda_1(x,y)$ is even in $y$.
Therefore $y\mapsto \lambda_1(x,y)$ is nondecreasing on $[0,\infty)$, and
since $s(x,y)=\sqrt{\lambda_1(x,y)}\ge0$, the same is true of $y\mapsto s(x,y)$.
\end{proof}

\begin{proof}[Proof of Theorem~\ref{thm:main}]
Combine Theorem~\ref{thm:vertical} with
Proposition~\ref{prop:topological-reduction}.
\end{proof}

\begin{lemma}[Eventual connectedness]\label{lem:eventual-connectedness}
For this lemma, write
\[
\Sigma_\varepsilon^{(n)}
=
\sigma_\varepsilon\bigl(\mathbf A_n(a)\bigr).
\]
If
\[
2\sqrt a\,\sin\!\left(\frac{\pi}{2(n+1)}\right)<\varepsilon,
\]
then \(\Sigma_\varepsilon^{(n)}\) is connected.  Consequently, the set
\[
\mathcal N_*(a,\varepsilon)
=
\left\{
n\ge2:\;
2\sqrt a\,\sin\!\left(\frac{\pi}{2(n+1)}\right)<\varepsilon
\right\}
\]
is nonempty, and, with
\[
N_*(a,\varepsilon):=\min \mathcal N_*(a,\varepsilon),
\]
one has
\[
n\ge N_*(a,\varepsilon)
\quad\Longrightarrow\quad
\sigma_\varepsilon\bigl(\mathbf A_n(a)\bigr)\ \text{is connected}.
\]
\end{lemma}

\begin{proof}
Fix \(n\ge2\), and write
\[
\mathbf A=\mathbf A_n(a).
\]
By Lemma~\ref{lem:real-spectrum}, the eigenvalues of \(\mathbf A\) are real, simple,
and given by
\[
\lambda_k^{(n)}=2\sqrt a\cos\!\left(\frac{k\pi}{n+1}\right),
\qquad
k=1,\dots,n,
\]
with
\[
\lambda_1^{(n)}>\lambda_2^{(n)}>\cdots>\lambda_n^{(n)}.
\]

For each \(k\in\{1,\dots,n\}\), choose a unit eigenvector \(\bm v_k\) of \(\mathbf A\)
corresponding to \(\lambda_k^{(n)}\).  Then, for every \(z\in\mathbb C\),
\[
s_{\min}(z\mathbf I_n-\mathbf A)
\le
\|(z\mathbf I_n-\mathbf A)\bm v_k\|_2
=
|z-\lambda_k^{(n)}|.
\]
Hence the open disk
\[
D_k:=\left\{z\in\mathbb C:\ |z-\lambda_k^{(n)}|<\varepsilon\right\}
\]
satisfies
\[
D_k\subset \Sigma_\varepsilon^{(n)},
\qquad k=1,\dots,n.
\]

Now compute the consecutive eigenvalue gaps.  Using
\[
\cos\alpha-\cos\beta
=
2\sin\!\left(\frac{\alpha+\beta}{2}\right)
 \sin\!\left(\frac{\beta-\alpha}{2}\right),
\]
we obtain, for \(k=1,\dots,n-1\),
\begin{align*}
\lambda_k^{(n)}-\lambda_{k+1}^{(n)}
&=
2\sqrt a
\left(
\cos\!\left(\frac{k\pi}{n+1}\right)
-
\cos\!\left(\frac{(k+1)\pi}{n+1}\right)
\right) \\
&=
4\sqrt a\,
\sin\!\left(\frac{(2k+1)\pi}{2(n+1)}\right)
\sin\!\left(\frac{\pi}{2(n+1)}\right) \\
&\le
4\sqrt a\,
\sin\!\left(\frac{\pi}{2(n+1)}\right).
\end{align*}
Therefore, if
\[
2\sqrt a\,\sin\!\left(\frac{\pi}{2(n+1)}\right)<\varepsilon,
\]
then
\[
\lambda_k^{(n)}-\lambda_{k+1}^{(n)}<2\varepsilon,
\qquad
k=1,\dots,n-1.
\]
Hence
\[
D_k\cap D_{k+1}\neq\varnothing,
\qquad
k=1,\dots,n-1,
\]
so the union
\[
\mathcal D:=\bigcup_{k=1}^n D_k
\]
is connected.

We next show that every connected component of \(\Sigma_\varepsilon^{(n)}\) contains
an eigenvalue of \(\mathbf A\).  First,
\[
s_{\min}(z\mathbf I_n-\mathbf A)\ge |z|-\|\mathbf A\|_2,
\]
and therefore \(\Sigma_\varepsilon^{(n)}\) is bounded.  Let \(\Omega\) be a connected
component of \(\Sigma_\varepsilon^{(n)}\), and suppose that
\[
\Omega\cap\sigma(\mathbf A)=\varnothing.
\]
Then
\[
\varphi(z):=\|(z\mathbf I_n-\mathbf A)^{-1}\|_2
\]
is continuous on \(\overline{\Omega}\), subharmonic on \(\Omega\), and satisfies
\[
\varphi(z)>\varepsilon^{-1},
\qquad
z\in\Omega.
\]
Since \(\Omega\) is a connected component of the open set \(\Sigma_\varepsilon^{(n)}\),
we have
\[
\partial\Omega\subset\partial\Sigma_\varepsilon^{(n)}.
\]
Because
\[
\Sigma_\varepsilon^{(n)}
=
\left\{
z\in\mathbb C:\ s_{\min}(z\mathbf I_n-\mathbf A)<\varepsilon
\right\},
\]
continuity of \(s_{\min}\) yields
\[
s_{\min}(z\mathbf I_n-\mathbf A)=\varepsilon,
\qquad
z\in\partial\Sigma_\varepsilon^{(n)},
\]
and hence
\[
\varphi(z)=\varepsilon^{-1},
\qquad
z\in\partial\Omega.
\]
This contradicts the maximum principle for subharmonic functions on bounded
domains.  Therefore every connected component of \(\Sigma_\varepsilon^{(n)}\)
contains an eigenvalue of \(\mathbf A\).

Since \(\mathcal D\subset\Sigma_\varepsilon^{(n)}\) is connected, it is contained in a single
connected component of \(\Sigma_\varepsilon^{(n)}\).  But \(\mathcal D\) contains all
eigenvalues \(\lambda_k^{(n)}\), so every connected component of
\(\Sigma_\varepsilon^{(n)}\) contains an eigenvalue lying in \(\mathcal D\), and therefore
must coincide with that one component.  Thus \(\Sigma_\varepsilon^{(n)}\) is connected.

Finally,
\[
\sin\!\left(\frac{\pi}{2(n+1)}\right)\longrightarrow0
\qquad\text{as }n\to\infty,
\]
so \(\mathcal N_*(a,\varepsilon)\) is nonempty.  The last assertion follows immediately
from the definition of \(N_*(a,\varepsilon)\).
\end{proof}

\begin{lemma}[Connectedness of the real slice]\label{lem:real-slice-connected}
Assume Theorem~\ref{thm:vertical}.  Define the real slice
\[
I_\varepsilon
:=
\Sigma_\varepsilon\cap\mathbb R
=
\left\{
x\in\mathbb R:\ s(x,0)<\varepsilon
\right\}.
\]
Then
\[
\Sigma_\varepsilon\ \text{is connected}
\quad\Longleftrightarrow\quad
I_\varepsilon\ \text{is connected}.
\]
\end{lemma}

\begin{proof}
We first show that
\[
\pi(\Sigma_\varepsilon)=I_\varepsilon,
\qquad
\pi:\mathbb C\to\mathbb R,\ \pi(x+iy)=x.
\]
Indeed, let \(z=x+iy\in\Sigma_\varepsilon\).  Then
\[
s(x,y)<\varepsilon.
\]
By Lemma~\ref{lem:conj},
\[
s(x,|y|)=s(x,y)<\varepsilon.
\]
Since Theorem~\ref{thm:vertical} states that \(y\mapsto s(x,y)\) is nondecreasing on
\([0,\infty)\), we obtain
\[
s(x,0)\le s(x,|y|)<\varepsilon,
\]
so \(x\in I_\varepsilon\).  This proves \(\pi(\Sigma_\varepsilon)\subset I_\varepsilon\).

Conversely, if \(x\in I_\varepsilon\), then
\[
s(x,0)<\varepsilon,
\]
hence \(x+0i\in\Sigma_\varepsilon\), and therefore \(x\in\pi(\Sigma_\varepsilon)\).
Thus \(\pi(\Sigma_\varepsilon)=I_\varepsilon\).

Assume now that \(\Sigma_\varepsilon\) is connected.  Since \(\pi\) is continuous,
\[
I_\varepsilon=\pi(\Sigma_\varepsilon)
\]
is connected.

Conversely, assume that \(I_\varepsilon\) is connected.  Since \(I_\varepsilon\subset\mathbb R\),
it is an interval.

Next we show that every vertical segment from a point of \(\Sigma_\varepsilon\) to the
real axis stays inside \(\Sigma_\varepsilon\).  Let \(z=x+iy\in\Sigma_\varepsilon\), and define
\[
\eta_{x,y}:[0,1]\to\mathbb C,
\qquad
\eta_{x,y}(\tau)=x+i\tau y.
\]
For \(\tau\in[0,1]\), Lemma~\ref{lem:conj} and Theorem~\ref{thm:vertical} give
\[
s(x,\tau y)=s(x,|\tau y|)
\le s(x,|y|)
= s(x,y)
<\varepsilon.
\]
Hence
\[
\eta_{x,y}([0,1])\subset\Sigma_\varepsilon.
\]

Take any two points
\[
z_0=x_0+iy_0\in\Sigma_\varepsilon,
\qquad
z_1=x_1+iy_1\in\Sigma_\varepsilon.
\]
By the first part of the proof,
\[
x_0,x_1\in I_\varepsilon.
\]
Since \(I_\varepsilon\) is an interval, the horizontal segment
\[
\gamma(\tau)=(1-\tau)x_0+\tau x_1,
\qquad
\tau\in[0,1],
\]
lies entirely in \(I_\varepsilon\subset\Sigma_\varepsilon\).

Now concatenate the three paths
\[
\tau\mapsto x_0+i(1-\tau)y_0,
\qquad
\tau\mapsto (1-\tau)x_0+\tau x_1,
\qquad
\tau\mapsto x_1+i\tau y_1.
\]
The first and third lie in \(\Sigma_\varepsilon\) by the vertical-segment property proved
above, and the middle one lies in \(\Sigma_\varepsilon\) because it is contained in
\(I_\varepsilon\).  Thus \(z_0\) and \(z_1\) can be joined by a path in \(\Sigma_\varepsilon\).

Therefore \(\Sigma_\varepsilon\) is path-connected, hence connected.
\end{proof}

\section{Pentadiagonal determinant formulas for \texorpdfstring{$(xI_n-A_n(a))^T(xI_n-A_n(a))-tI_n$}{(xI\_n-A\_n(a)) transpose (xI\_n-A\_n(a)) - tI\_n}}

Let
\[
A_n(a)=\operatorname{tridiag}(a,0,1)\in \mathbb{R}^{n\times n},
\qquad 0<a<1,
\]
and define
\[
K_n(x):=(xI_n-A_n(a))^T(xI_n-A_n(a)),
\qquad
M_n(x,t):=K_n(x)-tI_n.
\]
Set
\[
\rho:=x^2+a^2-t,\qquad
p:=x^2+a^2+1-t,\qquad
q:=-(1+a)x,\qquad
s:=a,\qquad
\widetilde\rho:=x^2+1-t.
\]
Then $M_n(x,t)$ is the symmetric pentadiagonal matrix
\[
M_n(x,t)=
\begin{pmatrix}
\rho & q & s \\
q & p & q & s \\
s & q & p & \ddots \\
& \ddots & \ddots & \ddots & q \\
& & s & q & \widetilde\rho
\end{pmatrix},
\]
with the obvious zero entries outside the bandwidth~$2$.

\medskip

Define the one-corner auxiliary matrix
\[
E_n(\rho):=
\begin{pmatrix}
\rho & q & s \\
q & p & q & s \\
s & q & p & \ddots \\
& \ddots & \ddots & \ddots & q \\
& & s & q & p
\end{pmatrix},
\qquad
e_n(\rho):=\det E_n(\rho),
\]
and let $u_j$ denote the $j$th Euclidean basis vector.

\begin{proposition}[Rank-one reduction to a one-corner family]
For every $n\ge 2$,
\[
M_n(x,t)=E_n(\rho)+(\widetilde\rho-p)\,u_nu_n^T
       =E_n(\rho)-a^2\,u_nu_n^T.
\]
Consequently,
\[
\det M_n(x,t)=e_n(\rho)-a^2 e_{n-1}(\rho).
\]
\end{proposition}

\begin{proof}
By inspection, $M_n(x,t)$ and $E_n(\rho)$ differ only in the $(n,n)$ entry, and
\[
\widetilde\rho-p=(x^2+1-t)-(x^2+a^2+1-t)=-a^2.
\]
Hence
\[
M_n(x,t)=E_n(\rho)-a^2u_nu_n^T.
\]
Since only the $(n,n)$ entry is changed, determinant multilinearity in the last row
(or, equivalently, cofactor expansion in the $(n,n)$ entry) gives
\[
\det M_n(x,t)=\det E_n(\rho)-a^2\,\operatorname{cof}_{nn}(E_n(\rho)).
\]
Removing the last row and last column from $E_n(\rho)$ leaves precisely
$E_{n-1}(\rho)$, hence
\[
\operatorname{cof}_{nn}(E_n(\rho))=\det E_{n-1}(\rho)=e_{n-1}(\rho),
\]
and the claim follows.
\end{proof}

\begin{proposition}[Five-term recurrence for the one-corner determinant]
Assume $q\neq 0$ (equivalently, $x\neq 0$). Then for $n\ge 7$,
\[
e_n(\rho)
=
(p-s)e_{n-1}(\rho)
+(ps-q^2)\bigl(e_{n-2}(\rho)-s\,e_{n-3}(\rho)\bigr)
+s^3(s-p)e_{n-4}(\rho)
+s^5e_{n-5}(\rho),
\]
with initial values
\begin{align*}
e_1(\rho)&=\rho,\\
e_2(\rho)&=p\rho-q^2,\\
e_3(\rho)&=p^2\rho-q^2(\rho-2s)-p(q^2+s^2),\\
e_4(\rho)&=p^3\rho-p^2(q^2+s^2)-p\!\left(2q^2(\rho-s)+\rho s^2\right)
          +q^4+2q^2s(\rho-s)+s^4,\\
e_5(\rho)&=p^4\rho+q^4(\rho-4s)+\rho s^4+2q^2s^2(-\rho+s)-p^3(q^2+s^2)\\
&\qquad\quad
          +p\!\left(2q^4+4q^2\rho s+s^4\right)
          +p^2\!\left(-2\rho s^2+q^2(-3\rho+2s)\right).
\end{align*}
\end{proposition}

\begin{proof}
This is the one-corner pentadiagonal determinant recurrence specialized to
\[
(r,p,q,s)=(\rho,\ x^2+a^2+1-t,\ -(1+a)x,\ a).
\]
\end{proof}

\begin{corollary}[Closed form for the shifted determinant: generic case]
Assume
\[
q\neq 0,
\qquad
q^2\notin\left\{4s(p-2s),\ \frac{(p+2s)^2}{4}\right\}.
\]
Define
\[
\alpha:=\sqrt{(p+2s)^2-4q^2},
\qquad
\beta_1:=\sqrt{2(p-2s)(p+2s-\alpha)-4q^2},
\]
\[
\beta_2:=\sqrt{2(p-2s)(p+2s+\alpha)-4q^2},
\]
and
\[
\mu_1:=\frac{p-2s-\alpha-\beta_1}{4},\qquad
\mu_2:=\frac{p-2s-\alpha+\beta_1}{4},
\]
\[
\mu_3:=\frac{p-2s+\alpha-\beta_2}{4},\qquad
\mu_4:=\frac{p-2s+\alpha+\beta_2}{4},\qquad
\mu_5:=s.
\]
Then the numbers $\mu_1,\dots,\mu_5$ are the distinct roots of the characteristic
polynomial
\[
\chi(z):=
z^5-(p-s)z^4-(ps-q^2)(z^3-sz^2)-s^3(s-p)z-s^5,
\]
and therefore
\[
e_n(\rho)=\sum_{j=1}^5 \kappa_j \mu_j^n,
\qquad n\ge 1,
\]
where the coefficients $\kappa_1,\dots,\kappa_5$ are uniquely determined by
\[
\sum_{j=1}^5 \kappa_j \mu_j^m=e_m(\rho),
\qquad m=1,2,3,4,5.
\]
Consequently,
\[
\det M_n(x,t)
=
\sum_{j=1}^5 \kappa_j(\mu_j-a^2)\mu_j^{\,n-1},
\qquad n\ge 1.
\]
\end{corollary}

\begin{proof}
Since the roots of $\chi$ are distinct in the stated generic regime, the standard
characteristic-roots ansatz for the linear recurrence yields
\[
e_n(\rho)=\sum_{j=1}^5 \kappa_j\mu_j^n.
\]
The coefficients are fixed uniquely by matching the first five values
$e_1(\rho),\dots,e_5(\rho)$. The determinant identity then follows from
\[
\det M_n(x,t)=e_n(\rho)-a^2e_{n-1}(\rho)
\]
proved in Proposition~1.
\end{proof}

\begin{proposition}[Degenerate line $x=0$]
If $x=0$, then $q=0$ and the recurrence simplifies to
\[
e_n(\rho)=p\,e_{n-1}(\rho)-p s^2 e_{n-3}(\rho)+s^4 e_{n-4}(\rho),
\qquad n\ge 6,
\]
with initial values
\[
e_1(\rho)=\rho,\qquad
e_2(\rho)=p\rho,\qquad
e_3(\rho)=p(p\rho-s^2),\qquad
e_4(\rho)=(p^2-s^2)(p\rho-s^2).
\]
If moreover $p^2\neq 4s^2$, then
\[
\nu_1=-s,\qquad \nu_2=s,\qquad
\nu_3=\frac{p-\sqrt{p^2-4s^2}}{2},\qquad
\nu_4=\frac{p+\sqrt{p^2-4s^2}}{2}
\]
are distinct and
\[
e_n(\rho)=\sum_{j=1}^4 \widetilde\kappa_j \nu_j^n,
\qquad n\ge 1,
\]
where the coefficients $\widetilde\kappa_1,\dots,\widetilde\kappa_4$ are determined by
\[
\sum_{j=1}^4 \widetilde\kappa_j \nu_j^m=e_m(\rho),
\qquad m=1,2,3,4.
\]
Hence
\[
\det M_n(0,t)
=
\sum_{j=1}^4 \widetilde\kappa_j(\nu_j-a^2)\nu_j^{\,n-1},
\qquad n\ge 1.
\]
\end{proposition}

\begin{proof}
This is the $q=0$ specialization of the one-corner pentadiagonal recurrence, together
with the standard characteristic-roots representation of its solutions. The determinant
formula again follows from Proposition~1.
\end{proof}

\begin{proposition}[Alternative symmetric-corner reduction]
Define the symmetric-corner auxiliary matrix
\[
\widetilde D_n(\widetilde\rho):=
\begin{pmatrix}
\widetilde\rho & q & s \\
q & p & q & s \\
s & q & p & \ddots \\
& \ddots & \ddots & \ddots & q \\
& & s & q & \widetilde\rho
\end{pmatrix},
\qquad
\widetilde d_n(x,t):=\det \widetilde D_n(\widetilde\rho).
\]
Then
\[
M_n(x,t)=\widetilde D_n(\widetilde\rho)-(1-a^2)u_1u_1^T.
\]
If $\widetilde D_n(\widetilde\rho)$ is invertible, then
\[
\det M_n(x,t)
=
\widetilde d_n(x,t)\Bigl(1-(1-a^2)m_n(x,t)\Bigr),
\]
where
\[
m_n(x,t):=u_1^T\widetilde D_n(\widetilde\rho)^{-1}u_1.
\]
\end{proposition}

\begin{proof}
Since the only difference between $M_n(x,t)$ and $\widetilde D_n(\widetilde\rho)$ is in the
$(1,1)$ entry,
\[
\rho-\widetilde\rho=(x^2+a^2-t)-(x^2+1-t)=-(1-a^2),
\]
hence
\[
M_n(x,t)=\widetilde D_n(\widetilde\rho)-(1-a^2)u_1u_1^T.
\]
Now apply the matrix determinant lemma:
\[
\det(\widetilde D_n+uv^T)=\det(\widetilde D_n)\bigl(1+v^T\widetilde D_n^{-1}u\bigr),
\]
with $u=-(1-a^2)u_1$ and $v=u_1$.
\end{proof}

\begin{remark}
In the nongeneric repeated-root cases
\[
q^2\in\left\{4s(p-2s),\ \frac{(p+2s)^2}{4}\right\}
\]
(or, in the line $x=0$, when $p^2=4s^2$), the same procedure yields the usual
polynomial-times-exponential solutions, i.e.\ terms of the form
$n^\ell \mu^n$ associated with repeated characteristic roots $\mu$.
\end{remark}

\begin{lemma}[Even-size central reduction]\label{lem:even-central-reduction}
Let \(n=2m\ge 2\), \(0<a<1\), and let
\[
A_{2m}(a)=\operatorname{tridiag}(a,0,1).
\]
Let \(J_m\in\mathbb R^{m\times m}\) be the nilpotent shift
\[
(J_m)_{j,j+1}=1,\qquad j=1,\dots,m-1,
\]
and define
\[
L_m:=aI_m+J_m,\qquad U_m:=I_m+aJ_m^{\mathsf T}.
\]
Let \(\Pi_{2m}\) be the odd-even permutation matrix corresponding to the reordered basis
\[
e_1,e_3,\dots,e_{2m-1},e_2,e_4,\dots,e_{2m}.
\]
Then
\[
\Pi_{2m}^{\mathsf T}A_{2m}(a)\Pi_{2m}
=
\begin{bmatrix}
0 & U_m\\
L_m & 0
\end{bmatrix}.
\]
Consequently,
\[
s_{2m}(0)=\sigma_{\min}\!\bigl(A_{2m}(a)\bigr)=\sigma_{\min}(L_m)=\sigma_{\min}(aI_m+J_m).
\]
Moreover,
\[
s_{2m+2}(0)\le s_{2m}(0).
\]
\end{lemma}

\begin{proof}
The block form is immediate from the action of \(A_{2m}(a)\) on odd and even
coordinates: odd indices are mapped only to even indices, and even indices only to
odd indices. A direct inspection shows that the odd-even reordered matrix has
upper-right block \(U_m=I_m+aJ_m^{\mathsf T}\) and lower-left block
\(L_m=aI_m+J_m\).

Therefore
\[
\Pi_{2m}^{\mathsf T}A_{2m}(a)^{\mathsf T}A_{2m}(a)\Pi_{2m}
=
\begin{bmatrix}
L_m^{\mathsf T}L_m & 0\\
0 & U_m^{\mathsf T}U_m
\end{bmatrix},
\]
so
\[
\sigma_{\min}\!\bigl(A_{2m}(a)\bigr)
=
\min\{\sigma_{\min}(L_m),\sigma_{\min}(U_m)\}.
\]

It remains to compare \(U_m\) and \(L_m\). Take any \(x\in\mathbb C^m\), and define
\(y\in\mathbb C^m\) by
\[
y_j:=x_{m+1-j},\qquad j=1,\dots,m.
\]
Then \(\|y\|_2=\|x\|_2\). With the convention \(y_{m+1}:=0\), one has
\[
\|U_mx\|_2^2=\sum_{j=1}^m |y_j+a y_{j+1}|^2,
\qquad
\|L_my\|_2^2=\sum_{j=1}^m |a y_j+y_{j+1}|^2.
\]
Subtracting,
\[
\|U_mx\|_2^2-\|L_my\|_2^2
=
\sum_{j=1}^m\Bigl(|y_j+a y_{j+1}|^2-|a y_j+y_{j+1}|^2\Bigr)
=
(1-a^2)|y_1|^2\ge 0.
\]
Hence \(\|U_mx\|_2\ge \|L_my\|_2\) for every \(x\), and since \(x\mapsto y\) is an
isometry of \(\mathbb C^m\), it follows that
\[
\sigma_{\min}(U_m)\ge \sigma_{\min}(L_m).
\]
Thus
\[
s_{2m}(0)=\sigma_{\min}\!\bigl(A_{2m}(a)\bigr)=\sigma_{\min}(L_m).
\]

For the monotonicity in \(m\), note that \(L_m^{\mathsf T}L_m\) is the leading principal
\(m\times m\) submatrix of \(L_{m+1}^{\mathsf T}L_{m+1}\). By Cauchy interlacing,
\[
\lambda_{\min}\!\bigl(L_{m+1}^{\mathsf T}L_{m+1}\bigr)
\le
\lambda_{\min}\!\bigl(L_m^{\mathsf T}L_m\bigr).
\]
Taking square roots gives
\[
\sigma_{\min}(L_{m+1})\le \sigma_{\min}(L_m),
\]
and therefore
\[
s_{2m+2}(0)\le s_{2m}(0).
\qedhere
\]
\end{proof}

\end{document}
